% article-sample.tex — a complete LaTeX article, self-contained.
% Compile with: pdflatex article-sample.tex
%
% Everything below is invented for sample purposes. The institute, the survey,
% and the figures in it describe nothing real.

\documentclass[11pt,a4paper]{article}

\usepackage[utf8]{inputenc}
\usepackage[T1]{fontenc}
\usepackage{lmodern}
\usepackage{geometry}
\usepackage{hyperref}

\geometry{margin=2.5cm}

\title{Reading Habits in Municipal Libraries:\\
       A Two-Year Observational Survey}
\author{R.\ Almeida\thanks{Corresponding author, Fenwick Institute for Civic
        Research.} \and T.\ Okonkwo \and M.\ Halvorsen}
\date{14 March 2027}

\begin{document}

\maketitle

\begin{abstract}
We report on a two-year survey of borrowing behaviour across eleven municipal
library branches in the fictional district of Fenwick. Circulation records,
reservation queues, and door counts were collected monthly. The data suggest
that branch opening hours explain more of the variation in borrowing than
collection size does, and that evening hours are disproportionately valuable to
borrowers who visit fewer than four times a year. All figures in this paper are
fabricated for demonstration and must not be cited.
\end{abstract}

\tableofcontents

\section{Introduction}
\label{sec:intro}

A library branch is measured, most often, by the size of its collection. That is
a convenient measure because it is easy to count, but it is a poor predictor of
whether anybody walks through the door. Over two years we tracked eleven
branches that vary widely in both collection size and opening hours, and asked a
narrow question: which of the two moves circulation more?

The remainder of this paper is organised as follows. Section~\ref{sec:method}
describes how the records were gathered. Section~\ref{sec:results} presents the
observed relationships. Section~\ref{sec:discussion} argues that the result is
less about libraries than about scheduling, and Section~\ref{sec:conclusion}
notes what a follow-up study would need to measure.

\section{Method}
\label{sec:method}

\subsection{Sites}

Eleven branches were included. Three are open seven days a week, five are open
six days, and three close on both Sunday and Monday. Collection sizes range from
just under 9{,}000 volumes at the smallest branch to a little over 140{,}000 at
the district's central branch.

\subsection{Instruments}

Three instruments were used, in decreasing order of reliability:

\begin{enumerate}
  \item \textbf{Circulation records.} Every loan, renewal, and return is logged
        by the catalogue system with a timestamp accurate to the minute.
  \item \textbf{Reservation queues.} A snapshot of every outstanding hold was
        taken at 03:00 on the first of each month.
  \item \textbf{Door counts.} Infrared counters at each entrance, reset nightly.
        These over-count, because a visitor who steps outside and returns is
        counted twice.
\end{enumerate}

Door counts were therefore used only as a coarse cross-check, never as a
dependent variable.\footnote{An earlier draft of this paper did use them
directly. The reviewers were right to object.}

\subsection{Exclusions}

Two months were excluded from every branch: the December in which the catalogue
system was migrated, and the following February, when three branches ran a
fine-forgiveness campaign that distorted returns. That leaves twenty-two
usable months.

\section{Results}
\label{sec:results}

\subsection{Opening hours dominate}

Weekly opening hours correlate with annual loans per capita at
$r = 0.71$. Collection size correlates at $r = 0.24$. When both are entered into
a single model, collection size stops being distinguishable from noise, while
opening hours retains almost all of its explanatory power:

\begin{quote}
Each additional weekly opening hour is associated with roughly 340 additional
annual loans at a median-sized branch, holding collection size fixed.
\end{quote}

\subsection{Evening hours are not like daytime hours}

Splitting opening hours by time of day changes the picture. Hours after 18:00
are associated with about 2.4 times the loans of an equivalent hour between
10:00 and 16:00. The effect is concentrated in a group we labelled
\emph{occasional borrowers} --- those with fewer than four visits a year, who
account for 31\,\% of registered members and 9\,\% of loans.

\subsection{The smallest branch is the exception}

The Marlow Street branch has the smallest collection and the shortest hours, and
outperforms the model's prediction by a wide margin in every month observed. It
is also the only branch inside a building that houses something else --- a
community kitchen that runs four evenings a week. We do not treat this as a
finding, only as the most interesting thing we could not explain.

\section{Discussion}
\label{sec:discussion}

The result is unsurprising once stated plainly: a library that is shut cannot
lend anything. What makes it worth reporting is how weak the competing
explanation turned out to be. Collection size is the number that appears in
budget submissions, in press coverage, and in the branch comparisons published
by the district itself. It is the number that is easiest to grow with money and
hardest to grow with anything else.

Opening hours, by contrast, are a staffing question. They are unglamorous, they
recur every year, and they are the first thing cut when a budget tightens. Our
data suggest that this ordering is exactly backwards, at least where the goal is
circulation.

We are careful not to overstate this. Circulation is one purpose of a library
among several, and the branch that lends the most books is not automatically the
branch doing the most good. A study that took reference desk questions,
study-space occupancy, or programme attendance as its outcome might well order
the branches differently.

\section{Threats to validity}

\begin{itemize}
  \item \textbf{Self-selection.} Branches with longer hours may be sited where
        demand was already higher. We have no instrument that separates the two.
  \item \textbf{Two years is short.} Both years shared the same funding formula,
        so the range of variation in hours is narrower than it looks.
  \item \textbf{One district.} Fenwick has an unusual density of branches for
        its population, which may make the marginal branch less valuable here
        than it would be elsewhere.
\end{itemize}

\section{Conclusion}
\label{sec:conclusion}

Across eleven branches and twenty-two months, when a branch is open predicts
borrowing better than what the branch holds. Evening hours are worth more than
their count suggests, and they matter most to the borrowers who come least
often. A follow-up should extend the observation window past a funding cycle and
should treat the co-located community kitchen at Marlow Street as a hypothesis
rather than an anomaly.

\section*{Acknowledgements}

We thank the branch managers of the Fenwick district for two years of patient
data entry, and the anonymous reviewers who removed the door counts from
Section~\ref{sec:results}.

\begin{thebibliography}{9}

\bibitem{fenwick2025}
Fenwick District Council.
\newblock \emph{Annual Library Service Report}.
\newblock Fenwick, 2025.

\bibitem{okonkwo2026}
T.~Okonkwo.
\newblock Counting doors, counting people: instrumentation error in footfall
studies.
\newblock \emph{Journal of Invented Methods}, 12(3):211--229, 2026.

\bibitem{halvorsen2026}
M.~Halvorsen and R.~Almeida.
\newblock Opening hours as a policy lever.
\newblock In \emph{Proceedings of the Fictional Conference on Civic
Infrastructure}, pages 88--97, 2026.

\end{thebibliography}

\end{document}
